\subsection{The dyadic field contribution and characters of small order}
\label{adds:field-budget}

Fix a primitive target \(\psi\), write \(K\) for its abelian
field, and put \(t(v)=\lfloor v^{401}\rfloor\). For a prime \(p\)
and a number field \(L\) write
\[
 h_p(L,e)=\operatorname{length}_{\mathbf Z_p}
 S_p(L,(\mathbf Q_p/\mathbf Z_p)(-e)).
\]
The strict conditions here are those of Section~\ref{newcoh:section}.
They are invariant under changes of algebraic closure and prolongation.

\begin{theorem}[One window for the dyadic field and the fixed odd layers]
\label{adds:dyadic-window}
Fix \(0<\varepsilon<1/52\) and a nonnegative integer \(s\).
At finitely many fixed odd primes \(p\), prescribe abelian base fields
\(B_p/\mathbf Q\) containing \(\mu_p\), and fixed nonnegative layer
indices \(m_p\). Write \(B_{p,m_p}\) for the corresponding degree
\(p^{m_p}\) cyclotomic layer. There are \(J\ge1\) and \(X_0\),
depending only on these fixed choices and the target, with the
following property. For \(X\ge X_0\) put \(H=X^{1/J}\). Take any
distinct primes \(q\equiv3\bmod8,\ d\equiv7\bmod16\), coprime to
the target conductor, with \(\log(qd)\le H^{\varepsilon/4}\).
Allow any additional set \(\mathcal E_*\) of at most \(s\) positive
integers, chosen even after \(q,d\).

There is a \(W\ge H\), \(W^2\le X\), such that for every
\(v\in[W,W^2]\) and every page index \(r\) with
\(v^{2+8\varepsilon}\le2^r\le2v^{2+8\varepsilon}\), the actual field
\[
 E=K\mathbf Q(\mu_{2^{r+2}})^+(\sqrt{-q},\sqrt{-d})
\]
of Proposition~\ref{host:art} satisfies
\[
 \sum_{1\le e<t(v)} h_2(E,e)\le \exp(v^{4\varepsilon/5}),
 \qquad
 \sum_{1\le e<t(v)}h_p(B_{p,m_p},e)\le\exp(v^{4\varepsilon/5})
 \quad\text{for each prescribed }p .
\]
All these lengths are finite. On the same interval every
\(e\in\mathcal E_*\) with \(e\le v^{401}\) satisfies
\(e\le v^{\varepsilon/64}\).
\end{theorem}
\begin{proof}
Put \(F_0=K(i,\sqrt{-q},\sqrt{-d})\) and
\(F_\infty=F_0(\mu_{2^\infty})\). The moving-field argument in
Section~\ref{newcoh:section} bounds the number of \(\mathbf Z_2\)
generators of the split unramified module by
\(R_2=4(g_K+6[K:\mathbf Q])\), uniformly in \(q,d\).
Indeed the two primes have respectively \(2\) and \(4\) places in
\(\mathbf Q(\mu_{2^\infty})\), since
\(v_2(q^2-1)=3,\ v_2(d^2-1)=4\).
The fixed base change multiplies these numbers by at most
\([K:\mathbf Q]\). Each newly allowed inertia group contributes
at most one mod-two character, and passing to each quadratic
adjunction multiplies the dimension bound by at most two.
At each prescribed odd prime choose a generator bound \(R_p\ge1\)
for the split unramified module over \(B_p(\mu_{p^\infty})\).
These bounds follow from the fixed abelian \(\mu=0\) theorem and the
class-field argument in the same section~\cite{ferreroWashington}.

At any of these primes let \(j_p=\lfloor\log_p R_p\rfloor+3\).
For \(p=2\), the cyclotomic image over \(F_0\) is
\(1+2^c\mathbf Z_2\) for some \(c\ge2\), since \(i\in F_0\).
It follows that \(\mathbf Q(\mu_{2^c})\subset F_0\) and
\(2^{c-1}\le[F_0:\mathbf Q]\le8[K:\mathbf Q]\).
The layer \(F_m=F_0(\mu_{2^{c+m}})\) has degree \(2^m\) over \(F_0\).
In particular \(E(i)=F_m\) with \(m=\max(0,r+2-c)=O(1+\log v)\).
The field \(F_{j_2}\) has bounded degree and lies in the cyclotomic
field of modulus \(2^{c+j_2}qd\operatorname{cond}(K)\). The
cyclotomic discriminant bound and the tower discriminant formula give
\[
 \log(2|\operatorname{disc}F_{j_2}|)\ll_K 1+\log(qd)\ll_K v^{\varepsilon/4}
 \qquad(H\le v\le X).
\]
At the odd primes the layers \(B_{p,j_p}\) are fixed.

Set \(R=R_2+\sum_p R_p+s,\ \eta=\varepsilon/64\), where the sum
here ranges over the prescribed odd primes. Take \(J=4(802/\eta)^R\).
Apply Proposition~\ref{newcoh:centres} in each tower with
\(Z=\lfloor X^{401}\rfloor\). For each prime it gives at most \(R_p\)
exceptional positive indices in \([1,Z]\). Apply the finite-window
argument of Section~\ref{sup:section} to the union of all these
sets and \(\mathcal E_*\). With this \(J\) it gives one \(W\) as in
the statement for which every member of the union at most \(v^{401}\)
is at most \(v^\eta\).

By Lemma~\ref{newcoh:desc} the tower lengths are exactly the
coinvariant lengths in Proposition~\ref{newcoh:centres}.
A nonexceptional \(e<t(v)\) at its prime consequently costs
\(O(1+\log v)\) at the required layer, since
\(\log Z\le401J\log v\). The constants are uniform in the pair.
For an exceptional \(e<t(v)\), apply Theorem~\ref{newcoh:lemmaK}
at the low layer \(j_p\). It is totally imaginary and contains the
required roots of unity. By Poitou--Tate the dual of its strict
group embeds in the indicated \(H^2\). Its length is therefore at
most \(2\exp(Cv^{3\varepsilon/4})\): the degree is bounded,
\(e^{32}\le v^{\varepsilon/2}\), and the log discriminant costs
\(O(v^{\varepsilon/4})\). The growth inequality
\eqref{newcoh:growth} adds at most \(R_p(m-j_p)_+\) on passing to
the required layer, with \(m=m_p\) at the odd primes.

For descent from \(E(i)\) to the actual field \(E\), restriction
preserves strictness. If the extension is quadratic, its
restriction kernel on \(H^1\) is
\(H^1(C_2,A^{G_{E(i)}})\), where \(A=(\mathbf Q_2/\mathbf Z_2)(-e)\).
The invariants are finite cyclic since the cyclotomic image is open
and \(e>0\). Evaluation at the nonidentity element makes this
cohomology a subquotient of a cyclic group, and restriction followed
by corestriction shows it is killed by two. Thus the loss in length
is at most one per twist. There is zero loss if \(E(i)=E\).
Summing the preceding estimates for fewer than \(v^{401}\) twists
gives the claimed bounds, since \(3\varepsilon/4<4\varepsilon/5\).
\end{proof}

For the prime selection in Theorem~\ref{pr:pair}, first fix this
\(J\), then use initial scale \(H\) and interval exponent \(J\).
The selected pair has
\(\log(qd)\le H^{\varepsilon/5}(1+\log^2 H)\le H^{\varepsilon/4}\)
eventually. Its discarded-character bounds hold on \([H,X]\),
which contains the entire selected interval. The rescaling in
Section~\ref{sup:section} uses the same pair. In particular
the dyadic term in Section~\ref{famsel:complement} costs
\((\log 2)\exp(v^{4\varepsilon/5})\), and the fixed odd layers
may be chosen to be those for the exceptional primes of that section.
Corollary~\ref{adds:fixed-char} uses them for every Case IV character.
Proposition~\ref{addo:compare} controls passage to the field at odd primes.

\begin{corollary}[The dyadic restriction on a zeta host]
\label{adds:zeta-injection}
For the trivial target one can take \(R_2=24,\ j_2=7\).
Restriction of strict groups from \(E\) to \(E(i)\) is injective
at every positive twist.
\end{corollary}
\begin{proof}
Every \(\mathbf Q(\mu_{2^a})\), \(a\ge2\), has odd class number.
Start with the Gaussian integers. In a quadratic step \(L/K'\) of
this tower the unique dyadic prime is totally ramified. If \(L\) has
a nonzero unramified quadratic character then its space of such
characters has a nonzero invariant under the involution, since
\((\sigma-1)^2=0\) in characteristic two. Its fixed character defines
an unramified quadratic extension of \(L\) which is Galois of degree
four over \(K'\). This group is abelian and the inertia at the unique
dyadic prime has order two. Its fixed field gives an everywhere
unramified quadratic extension of \(K'\), contradicting the induction
hypothesis. The finite-layer descent in Section~\ref{newcoh:section}
then gives \(g_{\mathbf Q}=0\).

Globally \(\mathbf Q(\mu_{2^\infty})\cap
\mathbf Q(i,\sqrt{-q},\sqrt{-d})=\mathbf Q(i)\):
inertia at \(q\) and \(d\) flips the corresponding square root
independently, and the cyclotomic two-tower is unramified there.
Locally \(-q\equiv5\bmod8\) gives the unramified quadratic extension
of \(\mathbf Q_2\), while \(-d\equiv1\bmod8\) is a square.
Thus, at any dyadic place \(w\), the completion of \(E(i)\) is the
compositum of \(\mathbf Q_2(\mu_{2^{r+2}})\) and this unramified
quadratic field. The completion of \(E\) uses the subfield fixed
by inversion in the same cyclotomic field. Their local extension
has degree two. Since each higher cyclotomic layer is totally
ramified, the cyclotomic images of \(G_{E(i)}\) and
\(G_{E(i)_w}\) are both \(1+2^{r+2}\mathbf Z_2\).
The restriction kernels in global and local degree one are
therefore inflations from the very same \(C_2\) acting on the
same invariant subgroup of \(A\). Local inflation is injective.
A strict global class restricts to zero at the dyadic place,
so its restriction to \(E(i)\) detects the class.
\end{proof}

\begin{proposition}[An all-prime budget in terms of character order]
\label{adds:order}
On any sufficiently late page of Section~\ref{famsel:complement},
sum over distinct host characters \(\xi\) at each coefficient
prime \(\ell\), with \(1\le e<t(v)\). Write
\(a=\operatorname{ord}(\xi)\) and \(U=ea\). Include every entry
satisfying at least one of
\[
 \begin{gathered}
 U\le v^{\varepsilon/128},\qquad \ell U\le v^{\varepsilon/64},
 \qquad
 U\le v^{\varepsilon/10}\ \text{ and }\
                  \ell>\max(2a+1,4e+2).
 \end{gathered}
\]
For every finite choice of coefficient primes the sum of their
normalized strict weights \(w_{\ell,e,\xi}\) is at most
\(\exp(v^{4\varepsilon/5})\). The estimate is uniform in the
primitive conductors of the characters on that page.
\end{proposition}
\begin{proof}
For a fixed kernel of a host character let \(K_\xi\) be its cyclic
field of degree \(a=\operatorname{ord}(\xi)\). Set \(F=K_\xi(i)\),
\(n=[F:\mathbf Q]\le2a\), and \(M=4N\), so that
\(F\subset\mathbf Q(\mu_M)\). Grouping characters by kernel is
independent of the coefficient prime. For each of them, restriction
to \(F\), Poitou--Tate, and the even Quillen--Lichtenbaum comparison
in Proposition~\ref{newcoh:fullsupport} bound the normalized log
order of the image by \(\log\# K_{2e}(\mathcal O_F)\{\ell\}\).
The restriction kernel vanishes for \(\ell\nmid n\). Otherwise
put \(u_\ell=\max(1,v_\ell(M))+v_\ell(e)\).
Since \(F\subset\mathbf Q(\mu_M)\), its cyclotomic image contains
\(1+\ell^{\max(1,v_\ell(M))}\) as a unit (at two \(4\mid M\)).
The coefficient invariants over \(F\) have normalized log order
at most \(u_\ell\log\ell\). Inflation--restriction and the bound
by all functions on a group of order \(n\) cost at most
\(n u_\ell\log\ell\) for the kernel. Summing this cost uses
\(\sum_{\ell\mid n}u_\ell\log\ell\le\log(Mne)\).
At fixed character kernel there are
at most \(\varphi(a)\) characters at any \(\ell\). Summing yields
\[
 \sum_{\ell,\ \ker\xi\ {\rm fixed}} w_{\ell,e,\xi}
 \le\varphi(a)
       \bigl(\log\#K_{2e}(\mathcal O_F)+n\log(Mne)\bigr).
\]

For clarity, the degree-sensitive estimate for the \(K\)-group
follows from the two parts of the proof of Theorem~\ref{famsel:rate}.
Put \(j=2e,\ H_0=\max(n+1,2j+2)\). For \(\ell>H_0\),
\cite[Theorem 1.1]{bayerEmeryHouriet} bounds the sum of the
primary log orders by
\[
 B=(2j+1)^{71j^4 n^3}\,n^{293j^5 n^5}
          |\operatorname{disc}F|^{528j^5 n^4}.
\]
For \(\ell\le H_0\), adjoining \(\mu_\ell\) costs degree at
most \(H_0\) prime to \(\ell\) at that individual prime. Pullback
on the primary \(K\)-group is injective by the trace identity in
the same proof. The new field has degree at most \(nH_0\),
conductor at most \(MH_0\), and contains the required roots of
unity. The bound for the full primary \(K\)-group in the proof
of Theorem~\ref{newcoh:lemmaK} gives
\[
 \log\#K_{2e}(\mathcal O_F)
 \le B+H_0\exp\bigl((jnH_0)^{32}(1+nH_0)\log(2MH_0)\bigr).
\]
These bounds apply to the totally imaginary field \(F\) of
degree at least two. Write \(U=ea\ge1\). Then \(jn\le4U\),
\(H_0\le6U\), \(nH_0\le12U^2\) and
\(\log|\operatorname{disc}F|\le n\log M\).
Taking logarithms of the positive bounds on the right gives
\[
 \log\bigl(1+\log\#K_{2e}(\mathcal O_F)\bigr)
       \ll U^{66}\log(2MU)
\]
with an absolute constant. On a page, \(\log(Nt)\le2v^{\varepsilon/4}\).
Thus for \(U\le v^{\varepsilon/128}\) the last display is
\(O(v^{49\varepsilon/64})\). The sum of the multiplicities
\(\varphi(a)\) over the cyclic quotients that occur is at most
the number of host characters, at most
\(4[K:\mathbf Q]\,2^r\ll_\psi v^4\).
In the second range apply the same restriction bound separately
at \(\ell\) with \(F'=F(\mu_\ell)\). The exponent in the
small-prime estimate is now at most
\(C(\ell U)^{33}\log(2M\ell U)=O(v^{49\varepsilon/64})\):
here \(j[F':\mathbf Q]\le4\ell U\) and
\([F':\mathbf Q]\le2\ell U\). There are at most
\(v^{\varepsilon/64}\) primes to sum in this range. The restriction
kernel cost \(n\log(Mne)\) per character is also absorbed.
In the third range \(\ell>H_0\) and \(\ell\nmid n\), so the sum over
these primes costs at most \(\varphi(a)B\) for the kernel orbit.
In every factor defining \(B\), including the discriminant term,
\(\log B\ll U^5\log(2MU)=O(v^{3\varepsilon/4})\).
Multiplication by the polynomial counts of characters, twists and,
in the second range, primes is absorbed in \(\exp(v^{4\varepsilon/5})\),
since \(49/64<4/5\).
\end{proof}

At the zeta target this includes every quadratic host character
at every coefficient prime for \(2e\le v^{\varepsilon/128}\).
In particular the bound applies to \(J_d\) at three with the
moving discriminant \(-d\) independently of the exponent of
\(\operatorname{Cl}(\mathbf Q(\sqrt{3d}))[3^\infty]\). For
Case~V it includes every character of order at most
\(v^{\varepsilon/128}\), and also those of order at most
\(v^{\varepsilon/64}/\ell\), since \(e=1\) there. Entries outside
the three ranges still require the other estimates of
Section~\ref{famsel:complement}.

\begin{corollary}[Inducing fields on the quartic zeta tests]
\label{adds:quartic}
On a sufficiently late zeta host let \((\ell,e,\xi)\) be in Case II
and consider an absolutely irreducible point of the good quartic
problem in equal characteristic. If it becomes reducible over a
finite extension, then \(\xi\) is quadratic,
\(\ell-1\mid 2e\), and \(\ell-1\nmid e\). Writing \(\Delta_\xi\)
for the fundamental discriminant of \(\xi\), the point is induced
from a strongly irreducible plane over the real quadratic field
of discriminant
\[
            D=\Delta_\xi\,(-1)^{(\ell-1)/2}\ell .
\]
This same field is the base of the strict Selmer calculation:
Proposition~\ref{famsel:quadraticlayers} computes the initial layers
there, and Proposition~\ref{famsel:alltower} computes deeper layers
over its indicated cyclotomic extensions.
The number of possible input triples with \(2\le e<t\) is
\(O(t\log\log(3t))\), uniformly in the host. Among them the total
weight of those with
\(\exp\operatorname{Cl}(\mathbf Q(\sqrt D))[\ell^\infty]\mid
\ell^{v_\ell(e)}\) is \(O(t\log\log(3t)\log(2Nt))\).
\end{corollary}
\begin{proof}
The Case II primes are odd. The fixed exceptional ones are \(3,7\)
and the others satisfy \(\ell-1\mid2e\). Hence none equals the two
odd level primes \(q,d>\max(2t+1,7)\). The factor of the residual
character \(b=\bar\varepsilon_\ell^{\,e}\bar\xi\) at \(\ell\)
and its factor of level \(N\) thus have independent images by the
Chinese remainder theorem. Proposition~\ref{hr:case3strong} forces
\(b\) to be quadratic and ramified at \(\ell\). Its cyclotomic factor
therefore has order exactly two, giving both divisibility claims.
The host character has two-power order, which reduction at the odd
prime preserves. Case II requires its order to be even. It must
therefore be quadratic as well. The two primitive discriminants
are coprime, and their product is positive by the parity of the
input. This gives the stated \(D\) and identifies the field with
the first field in Proposition~\ref{famsel:quadraticlayers}.

The quadratic intertwiner splits into two planes over this field,
as in Proposition~\ref{hr:ind}. Each has residual semisimplification
a sum of characters. Suppose one plane \(U\) becomes reducible on
an open subgroup. On an open normal subgroup of \(G_{\mathbf Q(\sqrt D)}\)
contained in it, the restriction is semisimple: the translates of
an invariant line span the irreducible plane. It is thus a sum of
two characters. At every finite place their inertial images are
finite by the equal-characteristic inertia conditions in the proof
of Proposition~\ref{hr:ind}. Only finitely many places ramify.
A common power of each character is everywhere unramified and
factors through a finite narrow class group. Both characters have
finite image. Passing to powers in the finite coset quotient now
makes all eigenvalues of \(U(g)\) roots of unity for every \(g\).
In equal characteristic they are constants and reduce injectively.
The characteristic polynomials consequently agree with those of
the constant lifts of the two residual characters. Brauer--Nesbitt
contradicts the irreducibility of \(U\).

There are at most four quadratic host characters of each parity,
as in Corollary~\ref{famsel:zetasubsum}. For a fixed \(\ell\) at
most \(2t/(\ell-1)\) positive twists qualify.
For \(y=2t+1\) set \(s=1+1/\log y\). The Euler product gives
\(\sum_{p\le y}p^{-1}\le\mathrm e\log\zeta(s)
\le\mathrm e\log(1+\log y)\). This proves the claimed count.
Under the exponent condition Proposition~\ref{famsel:quadraticlayers}
gives weight \(\log|\operatorname{Cl}(\mathbf Q(\sqrt D))[\ell^\infty]|\),
at most \(\log(2Nt)\) by the ideal bound of
Corollary~\ref{famsel:zetasubsum}. Summing proves the last assertion.
\end{proof}

\begin{corollary}[The full Case IV character sum]\label{adds:fixed-char}
Let the primitive target have conductor \(M\), and fix any finite set
\(P\) of odd coefficient primes containing its exceptional primes
\(B_\psi\). The interval in Theorem~\ref{adds:dyadic-window} can
be chosen so that the sum of \(w_{\ell,e,\xi}\) over
\(\ell\in P\), \(1\le e<t(v)\), and all host characters of primitive
conductor dividing \(M\), is \(O_{\psi,P}(\exp(v^{4\varepsilon/5}))\).
In particular this bounds the entire Case IV contribution.
\end{corollary}
\begin{proof}
For each \(\ell\in P\), prescribe the fixed field
\(B_\ell=\mathbf Q(\mu_{4M\ell})\) and its base layer \(m_\ell=0\)
in Theorem~\ref{adds:dyadic-window}. A character of conductor dividing
\(M\) is trivial over \(B_\ell\). Restriction takes its strict group
into the strict group there after scalar extension. At the places of
the host level the class restricts to zero, and unramified classes
at the other primes away from \(\ell\) stay unramified. The kernel
bound in the proof of Proposition~\ref{newcoh:fullsupport}, applied
to this fixed field, gives
\[
 w_{\ell,e,\xi}\ \le\ (\log\ell)h_\ell(B_\ell,e)
                   +O_{\psi,P}(1+\log e).
\]
Indeed the group of the extension has bounded order, and the
cyclotomic image contains \(1+\ell^{v_\ell(4M\ell)}\). Its action
kills the coefficient invariants by
\(\ell^{v_\ell(4M\ell)+v_\ell(e)}\), so counting cocycles of the
finite quotient gives the asserted normalized kernel loss.
There are at most \(\varphi(M)\) distinct lifts at each prime.
Summing on the jointly selected interval proves the bound.

On a late host the primes \(q,d\) lie outside \(B_\psi\) and are
prime to \(2M\). Their local character factors have order at most
two by Proposition~\ref{host:art}; the local factor at two has
two-power order. Taking the prime-to-\(\ell\) part for odd \(\ell\)
preserves each such factor. Removing the local factor at
\(\ell\in B_\psi\) leaves them all in \(\chi_P\), which by Case IV
has odd order. They are consequently trivial. At each remaining
prime the conductor exponent is at most that of \(M\), since
the host conductor divides \(\operatorname{lcm}(2^{r+2},M)qd\).
Every Case IV character therefore occurs in the preceding sum.
\src{src/Development/AutomorphicArithmetic/AutomorphicVarieties/HeightAndVanishing/DirichletPrimeToLPartMulLPowerPart.lean}{86}{removing the local factor}
and \src{src/Development/AutomorphicArithmetic/AutomorphicVarieties/HeightAndVanishing/DirichletPrimeToLPartMulLPowerPart.lean}{217}{the Case IV order condition}.
\end{proof}
